% Vietnamese translation for OI Public Library
% Source-PDF: IOI中国国家候选队论文/2007/day2/1.杨弋《Hash在信息学竞赛中的一类应用》.pdf
\documentclass[11pt]{article}
\usepackage{oipl-vn}

\newcommand{\Oh}{\mathcal{O}}
\newcommand{\Hash}{\operatorname{Hash}}
\newcommand{\xor}{\mathbin{\mathrm{xor}}}

\title{Một lớp ứng dụng của hàm băm trong thi tin học}
\author{Yang Yi}
\date{2007}

\begin{document}
\maketitle

\origtitle{A class of applications of hashing in informatics contests}
\sourcepath{IOI China candidate team papers / 2007 / day 2 / Yang Yi}

\section*{Mở đầu}

Bảng băm là một cấu trúc dữ liệu rất hiệu quả và có phạm vi ứng dụng rộng.
Nếu hàm băm được thiết kế hợp lý, trong tình huống lý tưởng mỗi lần truy vấn
chỉ tốn thời gian $\Oh(h/r)$, tức tỉ lệ với thương giữa dung lượng bảng băm
và dung lượng còn trống. Chỉ cần dung lượng bảng băm lớn hơn khoảng $1.5$ lần
lượng dữ liệu thực sự được dùng, thời gian truy vấn về cơ bản có thể xem là
hằng số $\Oh(1)$.

Với một bảng băm, hàm băm tốt đặc biệt quan trọng vì nó bảo đảm hiệu quả của
bảng băm. Đặc trưng dễ thấy nhất của một hàm băm tốt là: những đối tượng
khác nhau sau khi băm chỉ có xác suất rất nhỏ cho cùng một giá trị. Nhờ đó ta
tránh được quá nhiều xung đột.

Bảng băm không thể tách khỏi hàm băm. Nhưng chiều ngược lại thì sao? Đôi khi
hàm băm lại có thể được dùng độc lập với bảng băm. Một ví dụ thường gặp là
thuật toán MD5 nổi tiếng: đó là một hàm băm, nhưng thường được dùng để mã hóa
tóm tắt thông tin và kiểm tra tính đúng đắn của thông tin. Nói cách khác, ta
trực tiếp so sánh kết quả do MD5 tính ra để suy đoán nội dung gốc có giống
nhau hay không. Ngoài MD5, mã kiểm tra CRC32 và thuật toán SHA-1 thông dụng
cũng dựa trên tư tưởng tương tự.

Vậy trong thi tin học, kiểu thuật toán này có đất dụng võ hay không?

Bài viết này bàn về loại hàm băm dùng với mục tiêu phát hiện trùng lặp hoặc
phán định tương đương. Ta bắt đầu bằng Ví dụ 1.

\section{Ví dụ 1: Ghép mẫu nhiều chiều}

\subsection{Mô tả bài toán}

Tìm vị trí xuất hiện đầu tiên của một xâu trong một xâu khác là việc rất đơn
giản: dùng KMP là đủ. Mở rộng lên hai chiều, bài toán trở thành tìm vị trí
xuất hiện đầu tiên của một ma trận trong một ma trận khác. Nếu mở rộng lên
$k$ chiều thì phải làm thế nào?

Mảng văn bản $X$ có kích thước theo các chiều là
$N_1,N_2,\ldots,N_k$, mảng mẫu $Y$ có kích thước theo các chiều là
$M_1,M_2,\ldots,M_k$. Đặt
\[
  N=N_1N_2\cdots N_k,\qquad M=M_1M_2\cdots M_k.
\]
Bảo đảm $k\le 10$, $N_i\ge M_i$, và cả $N$ lẫn $M$ đều không vượt quá
$500000$.

\subsection{Phân tích thuật toán}

Thuật toán thường gặp cho bài này là KMP nhiều chiều: trước hết tính tình
trạng khớp ở một chiều, rồi xử lý hai chiều, ba chiều, \ldots, cho đến $k$
chiều. Độ phức tạp thời gian là $\Oh(k(N+M))$. Tuy nhiên thuật toán này khó
hiểu, khó nhớ, và không dễ hoàn thành cũng như viết đúng trong vài giờ thi.

Cách làm vét cạn tương đối dễ cài đặt, vì thế cũng rất dễ nghĩ ra; dữ liệu
đối kháng cho nó lại rất dễ tạo. Do đó nó chỉ đáng thử khi hoàn toàn không
có ý tưởng nào khác.

Có lựa chọn thứ ba hay không? Câu trả lời là có. Cách vét cạn tương đương với
việc liệt kê mọi điểm bắt đầu, hiển nhiên số điểm bắt đầu không vượt quá $N$,
rồi kiểm tra từng điểm. Nhưng kiểm tra hai mảng con có giống nhau hay không
tốn tới $\Oh(M)$ thời gian. Đây chính là nguyên nhân khiến thuật toán vét cạn
rất chậm trên dữ liệu đối kháng. Ta có thể loại nhanh một số trường hợp hiển
nhiên không thể khớp trước khi so sánh hai mảng con hay không? Vì rất dễ tạo
dữ liệu chỉ khác một ký tự, bất kể so sánh theo thứ tự nào cũng có thể tốn
rất nhiều thời gian.

Hàm băm có thể phát huy tác dụng ở đây. Ta tính một giá trị băm cho mỗi mảng
con có cùng kích thước với mảng mẫu. Rõ ràng chỉ khi giá trị băm của một mảng
con bằng giá trị băm của mảng mẫu thì chúng mới đáng được so sánh.

KMP nhiều chiều bắt đầu từ trường hợp một chiều rồi mở rộng lên chiều cao
hơn. Với thuật toán ghép mẫu bằng hàm băm, ta cũng có thể bắt đầu từ một
chiều: đó chính là thuật toán Rabin--Karp.

Với một xâu $S$, có thể dùng hàm sau để tính giá trị băm:
\[
  f(S)=\sum_{i=1}^{\operatorname{len}(S)}
    S[i]\,p^{\operatorname{len}(S)-i}\pmod q.
\]
Giá trị băm của xâu mẫu $Y$ vì thế được tính rất dễ. Ký hiệu $S_i$ là xâu con
độ dài $M$ của xâu văn bản $X$ bắt đầu từ ký tự thứ $i$. Không khó thấy quan
hệ giữa $f(S_{i+1})$ và $f(S_i)$:
\[
  f(S_i)=\sum_{j=i}^{i+M-1} X[j]p^{i+M-1-j}\pmod q,
\]
\[
\begin{aligned}
  f(S_{i+1})
  &=\sum_{j=i+1}^{i+M} X[j]p^{i+M-j}\pmod q\\
  &=\left[
      p\sum_{j=i}^{i+M-1} X[j]p^{i+M-1-j}
      +X[i+M]-p^M X[i]
    \right]\pmod q\\
  &=\bigl[p f(S_i)+X[i+M]-p^M X[i]\bigr]\pmod q.
\end{aligned}
\]
Nhìn từ góc độ khác, nếu bỏ qua phép lấy dư theo $q$, hàm này xem xâu như
một số trong hệ cơ số $p$. Khi đó $S_{i+1}$ thu được bằng cách ``dịch trái''
$S_i$ một vị trí, bỏ ký tự đầu tiên rồi thêm ký tự mới ở bên phải. Vì vậy
công thức truy hồi trên cũng là hiển nhiên.

Nhờ công thức truy hồi này, ta có thể tính mọi giá trị băm của các xâu con
độ dài $M$ trong $X$ với thời gian tuyến tính.

Bây giờ mở rộng bài toán lên nhiều chiều. Số giá trị băm cần tính vẫn không
vượt quá $N$; vì thế chỉ cần có phương pháp truy hồi thích hợp, ta vẫn có thể
tính mọi giá trị băm của các mảng con trong thời gian tuyến tính. Thật ra,
một ma trận con $M_1$ hàng, $M_2$ cột có thể xem là một xâu gồm $M_2$ ``dải
dọc''. Ta trước hết tính giá trị băm cho mọi dải dọc độ dài $M_1$, rồi dùng
các giá trị đó để tính giá trị băm hai chiều.

Một điểm cần chú ý là giá trị $p$ dùng để tính băm một chiều, tức băm của các
dải dọc, và giá trị $p$ dùng để tính băm hai chiều không được giống nhau. Nếu
không, rất dễ nghĩ ra phản ví dụ sau:
\[
\begin{bmatrix}
a & b\\
a & a
\end{bmatrix}
\qquad\text{và}\qquad
\begin{bmatrix}
a & a\\
b & a
\end{bmatrix}.
\]
Hai ma trận này không giống nhau, nhưng kết quả băm chắc chắn bằng nhau nếu
hai chiều dùng cùng cơ số. Điều đó trái với mục đích dùng băm. Vì vậy ở chiều
thứ nhất và chiều thứ hai phải dùng các giá trị $p$ khác nhau.

Trong trường hợp hai chiều, tính mọi dải dọc độ dài $M_1$ tốn $\Oh(N)$ thời
gian; sau đó xem các giá trị băm của dải dọc như những ``chữ cái'' để tính
băm theo chiều ngang, tức giá trị băm của các ma trận con, cũng tốn $\Oh(N)$
thời gian. Tổng độ phức tạp vẫn là $\Oh(N)$.

Hàm băm hai chiều là, với $\operatorname{len}_1(S)$ và
$\operatorname{len}_2(S)$ lần lượt là kích thước của $S$ theo hai chiều:
\[
  f_2(S)=
  \sum_{i_1=1}^{\operatorname{len}_1(S)}
  \sum_{i_2=1}^{\operatorname{len}_2(S)}
  p_1^{\operatorname{len}_1(S)-i_1}
  p_2^{\operatorname{len}_2(S)-i_2}
  S[i_1,i_2]\pmod q.
\]
Ký hiệu $X_{i_1,i_2}$ là ma trận con của $X$ có góc trái trên tại hàng
$i_1$, cột $i_2$, gồm $M_1$ hàng và $M_2$ cột. Khi đó
\[
  f_2(X_{i_1,i_2})=
  \sum_{j_1=i_1}^{i_1+M_1-1}
  \sum_{j_2=i_2}^{i_2+M_2-1}
  p_1^{i_1+M_1-1-j_1}
  p_2^{i_2+M_2-1-j_2}
  X[j_1,j_2]\pmod q.
\]
Trượt cửa sổ sang phải một cột, ta có
\[
\begin{aligned}
  f_2(X_{i_1,i_2+1})
  &=
  \Biggl[
    p_2 f_2(X_{i_1,i_2})
    + \sum_{j_1=i_1}^{i_1+M_1-1}
      p_1^{i_1+M_1-1-j_1} X[j_1,i_2+M_2] \\
  &\qquad
    -p_2^{M_2}
      \sum_{j_1=i_1}^{i_1+M_1-1}
      p_1^{i_1+M_1-1-j_1} X[j_1,i_2]
  \Biggr]\pmod q.
\end{aligned}
\]
Hai tổng trong công thức lần lượt là giá trị băm một chiều của dải dọc mới
đi vào và dải dọc cũ bị bỏ ra, vì thế chúng có thể được tính trước.

Mở rộng lên $k$ chiều cũng tự nhiên như vậy. Trước hết tính mọi giá trị băm
của mảng con kích thước $M_1\times M_2\times\cdots\times M_{k-1}$, rồi xem
những giá trị băm ấy như các ký tự và dùng băm Rabin--Karp để tính giá trị
băm cho từng mảng con kích thước
$M_1\times M_2\times\cdots\times M_k$. Sau $k$ lượt tính toán, ta thu được
mọi giá trị băm cần thiết, với độ phức tạp thời gian $\Oh(kN+M)$.

Hàm băm ở $k$ chiều có thể viết là
\[
\begin{aligned}
  f_k(X_{i_1,i_2,\ldots,i_k})
  &=
  \sum_{j_1=i_1}^{i_1+M_1-1}
  \sum_{j_2=i_2}^{i_2+M_2-1}
  \cdots
  \sum_{j_k=i_k}^{i_k+M_k-1}
  \left(
    \prod_{t=1}^{k}p_t^{i_t+M_t-1-j_t}
  \right)\\
  &\qquad\qquad\cdot
  X[j_1,j_2,\ldots,j_k]\pmod q.
\end{aligned}
\]
Khi trượt theo chiều thứ $k$, đặt
\[
  G(c)=
  \sum_{j_1=i_1}^{i_1+M_1-1}\cdots
  \sum_{j_{k-1}=i_{k-1}}^{i_{k-1}+M_{k-1}-1}
  \left(\prod_{t=1}^{k-1}p_t^{i_t+M_t-1-j_t}\right)
  X[j_1,\ldots,j_{k-1},c]\pmod q,
\]
tức $G(c)$ là giá trị băm của lát $(k-1)$ chiều tại tọa độ $c$ theo chiều
cuối. Khi đó công thức truy hồi là
\[
  f_k(X_{i_1,\ldots,i_k+1})
  =
  \bigl[
    p_k f_k(X_{i_1,\ldots,i_k})
    +G(i_k+M_k)-p_k^{M_k}G(i_k)
  \bigr]\pmod q.
\]

Nếu không có thao tác lấy dư theo $q$, hàm băm này có thể được hiểu như một
phép chuyển cơ số. Khi mỗi $p_t$ đều lớn hơn kích thước bảng chữ cái tương
ứng, chỉ cần giá trị băm khác nhau là chắc chắn nội dung khác nhau. Nhưng
trên thực tế giá trị tính ra sẽ lớn đến mức làm thuật toán mất ý nghĩa: tốc
độ so sánh ``giá trị băm'' ấy không nhanh hơn so sánh trực tiếp. Vì vậy phép
lấy dư là cần thiết.

Trong tình huống lý tưởng, nếu một hàm băm có thể sinh ra $S$ giá trị khác
nhau thì xác suất hai mảng con khác nhau cho cùng một giá trị chỉ là $1/S$.
Hàm băm của ta dĩ nhiên chưa chắc đạt được lý tưởng ấy, nhưng từ biểu thức
$1/S$ ta dễ thấy rằng tăng $S$ có thể giảm hiệu quả xác suất phán đoán sai.
Từ đó còn rút ra một kết luận khác: nếu độ chính xác của một hàm băm chưa đủ,
chỉ cần tính thêm một hàm băm khác với nó là có thể nâng cao đáng kể độ đúng.
Dĩ nhiên trong bài này, độ phức tạp thời gian đã gần bằng
\[
  \Oh(kN+N\cdot (1/S)\cdot M),
\]
vì hàm băm không phải là lý tưởng tuyệt đối. Chỉ cần $S$ không nhỏ hơn đáng
kể so với $N/k$, hạng sau sẽ không phải nút thắt, và không cần bỏ quá nhiều
công sức vào hướng này. Ngược lại, tính nhiều hàm băm sẽ làm hằng số của
thuật toán tăng rõ rệt.

Về cách chọn $p$ và $q$, vì có thể hiểu quá trình này là chuyển sang cơ số
$p$ rồi lấy dư theo $q$, ta nên cố gắng tránh xung đột trong phạm vi cho
phép. Có thể dùng các gợi ý sau:
\begin{enumerate}
  \item $p$ không nên quá nhỏ, nếu không rất dễ sai.
  \item $q$ có thể chọn là một số nguyên tố, sau đó chọn $p$ sao cho với mọi
  $i$ từ $1$ đến $p-2$, ta có $p^i\bmod q\ne 1$.
\end{enumerate}

Hàm băm này không chỉ có thể được tính bằng cách lăn cửa sổ như trên, mà còn
có thể tính và duy trì bằng chia đoạn, chia khối, cây đoạn, v.v. Một số cách
cụ thể sẽ xuất hiện trong các ví dụ phía sau.

\section{Ví dụ 2: Equal squares (Ural 1486)}

\subsection{Mô tả bài toán}

Trong một ma trận ký tự $N\times M$, hãy tìm hai ma trận con hình vuông giống
nhau, được phép giao nhau, sao cho cạnh của hai hình vuông tìm được là lớn
nhất có thể.

\subsection{Phân tích thuật toán}

Sau kinh nghiệm ở Ví dụ 1, gặp bài này ta dễ nghĩ đến việc tìm kiếm nhị phân
độ dài cạnh hình vuông, rồi dùng bảng băm để kiểm tra độ dài ấy có khả thi
hay không. Hàm băm dùng giống trường hợp hai chiều ở Ví dụ 1. Độ phức tạp
thời gian là $\Oh(NM\log(NM))$.

Nhưng đáng tiếc là giới hạn thời gian của bài này rất chặt; chương trình làm
như vậy vẫn quá thời gian. Cuối cùng, sau khi dùng một sửa đổi có vẻ hơi mạo
hiểm, chương trình mới qua được: trực tiếp kiểm tra xem có xuất hiện hai giá
trị băm bằng nhau hay không. Nếu tồn tại hai giá trị băm bằng nhau thì coi
độ dài cạnh đó là khả thi; nếu không thì coi là không khả thi.

Bài này đã được thông qua như vậy, nhưng sửa đổi ấy có quá mạo hiểm không?
Thật ra, xác suất hai hình vuông con có nội dung khác nhau mà giá trị băm lại
bằng nhau là rất nhỏ. Vì thế trong phần lớn trường hợp ta có thể xem rằng
nếu hai giá trị băm bằng nhau thì nội dung gốc cũng bằng nhau.

Ở đây còn một vấn đề: bảng băm lưu cái gì? Nếu chỉ lưu kiểu Boolean mà không
nén bit thì có vẻ tốn khá nhiều bộ nhớ; hơn nữa bảng băm không thể mở quá
lớn, nên hai nội dung khác nhau vẫn có thể bị băm vào cùng một vị trí. Chỉ
so sánh ``vị trí này đã từng bị băm tới chưa'' vẫn chưa đủ an toàn. Nếu một
lần truy vấn có xác suất sai là một phần triệu, thì một triệu truy vấn đã tạo
ra xác suất sai đáng kể.

Một mẹo nhỏ là tính hai giá trị băm: một giá trị dùng để xác định vị trí
trong bảng băm, giá trị kia dùng để so sánh. Chỉ khi bắt đầu tìm ở một địa
chỉ trong bảng băm và tìm thấy một phần tử có giá trị băm so sánh bằng với
giá trị của mình, ta mới coi như phần tử ấy đã từng xuất hiện trong bảng.

\section{Ví dụ 3: Ghép mẫu không ổn định}

\subsection{Mô tả bài toán}

Cho hai xâu $A$ và $B$. Mỗi lần có thể thực hiện một trong các thao tác sau:
\begin{itemize}
  \item \texttt{INSERT i j}: chèn ký tự $j$ vào trước ký tự thứ $i$ của $A$.
  \item \texttt{DELETE i}: xóa ký tự thứ $i$ của $A$.
  \item \texttt{REVERSE i j}: đảo ngược đoạn từ vị trí $i$ đến vị trí $j$
  trong $A$.
  \item \texttt{QUERY i j}: hỏi độ dài lớn nhất có thể khớp giữa $A$ bắt đầu
  từ ký tự thứ $i$ và $B$ bắt đầu từ ký tự thứ $j$, tức độ dài LCP của hậu tố
  của $A$ bắt đầu tại $i$ và hậu tố của $B$ bắt đầu tại $j$.
\end{itemize}

\subsection{Phân tích thuật toán}

Nếu hai xâu không thay đổi, muốn tìm LCP chỉ cần xây dựng mảng hậu tố của
$A+B$. Nhưng trong bài này xâu $A$ liên tục thay đổi; hơn nữa từ các kiểu
thay đổi có thể thấy mỗi thao tác đều khiến mảng hậu tố biến đổi rất lớn. Vì
vậy ta nên tìm đường khác.

Với một $k$ cho trước, không khó kiểm tra LCP của hai xâu có dài ít nhất $k$
ký tự hay không: tính giá trị băm của $k$ ký tự đầu tiên của hai xâu rồi so
sánh chúng. Nếu bằng nhau, ta gần như có thể chắc chắn rằng LCP của hai xâu
dài ít nhất $k$; nếu khác nhau thì độ dài LCP chắc chắn nhỏ hơn $k$. Như vậy
có thể tính LCP bằng tìm kiếm nhị phân.

Vấn đề còn lại chuyển thành cách tính giá trị băm của một xâu con. Ta vẫn
dùng hàm băm Rabin--Karp. Nếu đã biết giá trị băm của $S_1$ và $S_2$, không
khó tính giá trị băm của $S_1+S_2$:
\[
\begin{aligned}
  f(S_1+S_2)
  &=
  \sum_{i=1}^{\operatorname{len}(S_1)+\operatorname{len}(S_2)}
  p^{\operatorname{len}(S_1)+\operatorname{len}(S_2)-i}
  (S_1+S_2)[i]\pmod q\\
  &=
  \left[
    \sum_{i=1}^{\operatorname{len}(S_1)}
    p^{\operatorname{len}(S_1)+\operatorname{len}(S_2)-i}S_1[i]
    +
    \sum_{i=1}^{\operatorname{len}(S_2)}
    p^{\operatorname{len}(S_2)-i}S_2[i]
  \right]\pmod q\\
  &=
  \bigl[
    p^{\operatorname{len}(S_2)}f(S_1)+f(S_2)
  \bigr]\pmod q.
\end{aligned}
\]
Các giá trị $p^t$ hiển nhiên có thể được tiền xử lý trong $\Oh(n)$ thời gian.
Vì vậy ta có thể dùng $\Oh(1)$ thời gian để ghép giá trị băm của hai xâu
thành giá trị băm của xâu nối.

Do đó, có thể dùng danh sách liên kết theo khối để duy trì giá trị băm của
các xâu con trong $A$. Cách làm tương tự như duy trì tổng từng phần, chỉ khác
ở chỗ giá trị băm xuôi và ngược của một xâu là khác nhau. Để có thể hoàn
thành thao tác \texttt{reverse} trong $\Oh(\sqrt n)$ thời gian, cần đảo ngược
một khối trong $\Oh(1)$ thời gian; vì thế mỗi khối phải duy trì hai giá trị
băm: một theo chiều xuôi và một theo chiều ngược. Ngoài điểm này, các thao
tác còn lại tương tự như duy trì tổng từng phần hoặc duy trì giá trị lớn
nhất. Khi đó thao tác chèn, xóa, đảo ngược đều tốn $\Oh(\sqrt n)$, còn truy
vấn tốn $\Oh(\sqrt n\log n)$.

Tương tự, một cách giải khác là duy trì giá trị băm trên cây splay. Mỗi khi
một nút được xoay hoặc cây con gốc tại nút ấy bị sửa đổi, ta tính lại giá trị
băm xuôi và băm ngược của nó. Nhờ vậy có thể đảo ngược một cây con trong
$\Oh(1)$ thời gian; thông qua thao tác tách, có thể đảo ngược một cây con
trong thời gian khấu hao $\Oh(\log n)$. Thao tác chèn và xóa hiển nhiên cũng
có thời gian khấu hao $\Oh(\log n)$, còn truy vấn có độ phức tạp khấu hao
$\Oh(\log^2 n)$.

\section{Ví dụ 4: Một lớp bài toán phán định đẳng cấu}

\subsection{Bài toán 1: so sánh hai cây có gốc}

Thuật toán dễ nghĩ là dùng hai xâu để biểu diễn hai cây. Nhưng nếu dùng băm
thì nên làm thế nào?

Ta có thể dùng một kiểu truy hồi trên cây để tính giá trị băm. Với một nút
$v$, trước hết tính giá trị băm của mọi con của nó, rồi sắp xếp các giá trị
đó tăng dần, ký hiệu là $H_1,H_2,\ldots,H_D$. Khi đó giá trị băm của $v$ có
thể được tính như sau:
\[
  \Hash(v)=
  (((\cdots(((a\cdot p\xor H_1)\bmod q)\cdot p\xor H_2)\bmod q)\cdots
  \cdot p\xor H_D)\cdot b)\bmod q.
\]
Nói cách khác, ta bắt đầu từ một hằng số, mỗi lần nhân với $p$, xor với một
phần tử, lấy dư theo $q$, rồi tiếp tục nhân với $p$, xor với phần tử kế tiếp,
lại lấy dư theo $q$, \ldots, cho đến phần tử cuối. Cuối cùng nhân kết quả với
$b$ rồi lấy dư theo $q$.

Lý do phải sắp xếp trước giá trị băm của các con là: chỉ khác thứ tự con
không làm hai cây khác nhau. Còn cách băm phía sau có thể tương đối tùy ý;
không nhất thiết phải dùng đúng hàm băm này, miễn là chọn một hàm băm đủ tốt.
Nhưng cần chú ý rằng hàm Rabin--Karp ở trên không thích hợp để dùng trực tiếp
tại đây, vì không khó tìm phản ví dụ.

Hãy xét hai cây sau bằng mô tả xâu ngoặc: cây trái có gốc với hai con, một
con là lá và một con có một lá; cây phải có gốc với hai con, một con có hai
lá và một con là lá. Nếu trong cây phải thứ tự con đúng như mô tả, thì vì
giá trị băm có thể xem là phân bố ngẫu nhiên nên có một nửa khả năng thứ tự
này xuất hiện sau khi sắp xếp. Giá trị băm của mọi lá bắt buộc giống nhau,
ký hiệu là $l$. Nếu dùng truy hồi
\[
  \Hash(v)=b\sum_{i=1}^{D}p^{D-i}v_i\pmod q,
\]
thì cây trái có giá trị
\[
  lb^2(p^2+p+1)(p+1)\pmod q,
\]
còn cây phải có giá trị
\[
  lb^2\bigl[(p+1)p+(p^3+p^2+p+1)\bigr]\pmod q.
\]
Hai giá trị này bằng nhau, nhưng hai cây lại không bằng nhau.

Tương tự, với cách tính ở phần đầu, nếu cuối cùng không nhân với $b$ thì cũng
sai; có thể phân tích trường hợp cây suy biến thành một đường thẳng.

Vậy, nếu so sánh trực tiếp giá trị băm chắc chắn có xác suất sai, vì sao còn
cần chỉ ra hàm băm nào dùng được, hàm nào không dùng được? Một số lỗi là
``sai số ngẫu nhiên'': không cần đổi phương pháp tính băm, chỉ đổi các hằng
như $p$ là có thể loại bỏ. Nhưng có những lỗi là ``sai số hệ thống'', do bản
thân hàm băm không hợp lý; trường hợp sau nên tránh tối đa. Dĩ nhiên trong
vài giờ thi không nhất thiết ta đánh giá được một hàm băm có hợp lý hay
không, nhưng tốt hơn nên chọn hàm băm mà khi chưa biết các hằng có thể thay
đổi như $p,q$, rất khó dựng phản ví dụ.

Bây giờ chỉ cần so sánh giá trị băm của hai cây là có thể phân biệt chúng có
giống nhau hay không. Độ phức tạp thời gian là $\Oh(n\log n)$.

Còn một hàm băm khác cũng giải được bài toán này, và rất hữu ích cho bài toán
tiếp theo. Ta xét việc biểu diễn cây bằng một xâu, tương tự biểu diễn nhỏ
nhất, rồi tính giá trị băm của xâu ấy. Tất nhiên nếu thật sự dựng toàn bộ xâu
như phương pháp biểu diễn nhỏ nhất thì hơi không đáng. Ta chỉ cần dựa vào giá
trị băm của các con để sắp xếp các con. Vì hàm băm Rabin--Karp xử lý được
giá trị băm của phép nối hai xâu, sau khi đã xác định thứ tự và biết số nút
của từng con, tức độ dài xâu con, không khó xác định giá trị băm của nút hiện
tại.

Cụ thể, nếu $c_1,c_2,\ldots,c_D$ là các con của nút hiện tại $v$, và $s(x)$
biểu diễn xâu chuyển đổi từ cây con gốc $x$, thì
\[
  s(v)=g(v)s(c_1)s(c_2)\cdots s(c_D).
\]
Ở đây $g(v)$ là một hàm nào đó có năng lực phân biệt nhất định về $v$, chẳng
hạn số con của $v$, hoặc một hàm băm nào đó liên quan tới $v$.

Với hàm băm Rabin--Karp
\[
  f(S)=\sum_{i=1}^{\operatorname{len}(S)}
  S[i]p^{\operatorname{len}(S)-i}\pmod q,
\]
ta có
\[
\begin{aligned}
  f(s(v))
  &=
  f(g(v)s(c_1)s(c_2)\cdots s(c_D))\\
  &=
  \left[
    p^{\operatorname{len}(s(v))-1}g(v)
    +
    \sum_{i=1}^{D}
      f(s(c_i))\,
      p^{\sum_{j=i+1}^{D}\operatorname{len}(s(c_j))}
  \right]\pmod q.
\end{aligned}
\]
Vì vậy chỉ cần một lần truy hồi trên cây là có thể tính giá trị băm.

\subsection{Bài toán 2: so sánh hai cây vô căn}

Lúc này nếu tiếp tục dùng biểu diễn nhỏ nhất thì độ phức tạp rất cao. Nhưng
nếu dùng băm hợp lý, độ phức tạp vẫn chỉ là $\Oh(n\log n)$.

Cách làm cụ thể như sau. Với mỗi cạnh $(a,b)\in E$, tính một giá trị băm
$F(a,b)$ sao cho nó phụ thuộc vào mọi $F(b,i)$ thỏa $i\ne a$ và $(b,i)\in E$.
Ở đây $F(a,b)$ và $F(b,a)$ là hai khái niệm khác nhau. Giá trị $F(a,b)$ dùng
để biểu diễn ``khi xem $a$ là cha của $b$, giá trị băm của cây con gốc
$b$''. Dễ thấy cách truy hồi này không gây chu trình và chắc chắn có thể tính
theo một thứ tự nào đó.

Trước hết chọn tùy ý một nút $r$ làm gốc. Với một cặp có thứ tự $(a,b)$ thỏa
$(a,b)\in E$, nếu $a$ gần gốc hơn $b$ thì gọi $(a,b)$ là hướng xuống, ngược
lại là hướng lên. Rõ ràng có thể dùng phương pháp thứ hai trong bài toán 1
để tính mọi $F(a,b)$ hướng xuống trong $\Oh(n\log n)$ thời gian.

Bây giờ xét gốc $r$. Mọi $(r,i)$ đều là hướng xuống nên mọi $F(r,i)$ đều đã
biết. Ta cần tính các giá trị $F(i,r)$. Với một con $t$ của $r$, giá trị
$F(t,r)$ có thể được tính từ giá trị băm của mọi con khác của $r$ ngoài $t$.
Nếu tính trực tiếp như vậy, trường hợp xấu nhất sẽ tốn $\Oh(n^2)$.

Ở đây dùng một mẹo nhỏ trong truy hồi trên cây. Vì mọi $F(r,i)$ đều đã biết,
ta sắp xếp chúng trước rồi xét. Hàm băm này về bản chất vẫn là băm của xâu;
ta sắp các điểm kề với $r$ theo $F(r,i)$, ký hiệu lần lượt là
$c_1,c_2,\ldots,c_D$, và ký hiệu $s(a,b)$ là xâu tương ứng với $F(a,b)$.

Có thể tìm kiếm nhị phân vị trí $k$ của $t$ trong
$c_1,c_2,\ldots,c_D$. Khi đó
\[
  s(t,r)=
  g(r)s(r,c_1)s(r,c_2)\cdots s(r,c_{k-1})
  s(r,c_{k+1})\cdots s(r,c_D).
\]
Các phần $g(r)$,
$s(r,c_1)s(r,c_2)\cdots s(r,c_{k-1})$ và
$s(r,c_{k+1})\cdots s(r,c_D)$ sau khi sắp xếp đều có thể được tiền xử lý
trong $\Oh(n)$ thời gian tương tự Ví dụ 1. Vì vậy có thể tính $F(t,r)$ trong
$\Oh(\log n)$ thời gian, do phải tìm kiếm nhị phân.

Khi một nút $v$ đã có $F(v,\operatorname{father}(v))$, thì tương đương với
việc mọi $F(v,i)$ với $i$ kề $v$ đều đã biết. Tình huống này giống hệt thảo
luận vừa rồi về gốc $r$, nên không cần lặp lại chi tiết.

Như vậy, lượt DFS thứ nhất xác định mọi $F(a,b)$ hướng xuống; lượt DFS thứ
hai xác định mọi $F(a,b)$ hướng lên. Độ phức tạp thời gian là bao nhiêu? Lượt
đầu hiển nhiên là $\Oh(n\log n)$. Ở lượt hai, với một nút có $D$ con, chi phí
là $\Oh(D\log D)<\Oh(D\log n)$; vì $\sum D=N$, tổng thời gian vẫn là
$\Oh(n\log n)$.

Giá trị băm của một đỉnh có thể được định nghĩa là giá trị băm của cả cây khi
lấy đỉnh ấy làm gốc. Sau khi đã tính mọi $F(a,b)$, giá trị băm của từng đỉnh
được tính không khó.

Cuối cùng, nếu hai cây vô căn có thể ghép từng đôi giá trị băm của đỉnh với
nhau, tức sau khi sắp xếp thì hoàn toàn tương ứng, hoặc kiểm tra bằng một
bảng băm, ta có thể coi chúng bằng nhau; nếu không thì coi là không bằng
nhau.

Vì giá trị băm của một đỉnh tương đương với việc lấy đỉnh ấy làm gốc, thu
được một cây có gốc rồi băm cây đó, có thể trực tiếp dùng giá trị băm đỉnh để
so sánh hai đỉnh có ở vị trí tương đương trong cây hay không. Cũng có thể
dựa trên sự tương ứng từng đỉnh để thu được ánh xạ giữa các đỉnh của hai cây
có cùng cấu trúc nhưng đánh số khác nhau. Do thay đổi một đỉnh sẽ làm giá trị
băm của toàn bộ các nút trên cây thay đổi, việc hai giá trị băm bằng nhau
cũng phản ánh rằng ``hai nút này rất có thể thuộc cùng một cây, hoặc thuộc
hai cây hoàn toàn giống nhau''. Kiểu băm này thậm chí có thể dùng để phán
định những điểm nào trong một rừng ở vị trí tương đương.

\subsection{Bài toán 3: so sánh hai đồ thị có hướng}

Vì trường hợp đồ thị vô hướng có thể chuyển thành đồ thị có hướng để xét, ở
đây chỉ bàn về đồ thị có hướng. Tác giả chưa tìm được một phương pháp băm
thật sự hiệu quả, nhưng có một thuật toán trong thực tế chưa gặp vấn đề.

Trước hết liệt kê một điểm bắt đầu $i$, đặt giá trị băm của mọi nút thành một
hằng số, chẳng hạn $1$. Sau đó lặp: mỗi lần giá trị băm mới của một nút được
tính là
\[
  f_{j+1}(v)=
  \left[
    a f_j(v)
    +b\sum_{\substack{w\in V\\ v\to w\in E}} f_j(w)
    +c\sum_{\substack{w\in V\\ w\to v\in E}} f_j(w)
    +g(v)
  \right]\bmod p,
\]
trong đó
\[
  g(v)=
  \begin{cases}
    d, & v=i,\\
    0, & v\ne i.
  \end{cases}
\]
Sau khi lặp $k$ lần, giá trị $f_k(i)$ thu được được xem là giá trị băm của
điểm $i$. Mấu chốt là phép tính $f_j(i)$ khác các giá trị khác, nhờ đó tránh
được nhiều trường hợp đặc biệt.

\paragraph{Bài liên quan: The labyrinth of wells (POI 1999/2000 Stage 2).}

Tóm tắt bài toán: có nhiều phòng; phòng trên cùng là điểm xuất phát, phòng
dưới cùng là điểm kết thúc. Trừ phòng dưới cùng, mỗi phòng có một trong ba
màu và có ba ống khác nhau dẫn tới các phòng bên dưới nó. Mọi phòng đều có
thể tới được, và mọi ống cuối cùng đều dẫn tới phòng dưới cùng. Nhưng một số
nút về bản chất là giống nhau: nếu chọn cùng một cách đi thì sẽ đi qua cùng
chuỗi màu. Hỏi sau khi gộp các nút bản chất giống nhau, số điểm còn lại ít
nhất là bao nhiêu.

Hai điểm có thể gộp khi và chỉ khi các đường đi phía sau chúng có màu hoàn
toàn giống nhau. Từ đó có thể thiết kế một hàm băm:
\[
  \Hash(x)=f(\Hash(x\to c_1),\Hash(x\to c_2),\Hash(x\to c_3),x\to color),
\]
trong đó $f$ là một hàm băm nào đó. Trong chương trình của tác giả, hàm băm
là
\[
  f(a,b,c,d)=((129(129a\xor b)\xor c)\xor 987654321)+d.
\]
Nhân tiện, để tăng hiệu suất chương trình, có thể dùng một số kỹ thuật thao
tác bit, chẳng hạn $a\cdot129=(a\ll 7)+a$.

\section{Ví dụ 5: Biểu thức tương đương (NOIP 2005)}

\subsection{Mô tả bài toán}

Cho một đa thức chưa rút gọn, rồi cho tiếp $n$ đa thức. Hãy xác định trong
$n$ đa thức ấy, những đa thức nào bằng đa thức ban đầu.

Ràng buộc: $n\le 26$, độ dài biểu thức không vượt quá $50$. Trong đa thức chỉ
có số, biến \texttt{a}, các phép \texttt{+}, \texttt{-}, \texttt{*},
\texttt{\^{}} (lũy thừa), dấu ngoặc trái phải và khoảng trắng. Các số đều
nhỏ hơn $10000$, số mũ trong phép lũy thừa không vượt quá $10$;
\texttt{+}, \texttt{-}, \texttt{*}, \texttt{\^{}}, \texttt{(}, \texttt{)}
đều chỉ đóng vai trò toán tử, và không có trường hợp lược bỏ dấu nhân.

\subsection{Phân tích thuật toán}

Cách làm thông thường là khai triển từng đa thức, gộp các hạng tử đồng dạng,
rồi so sánh. Thuật toán này tất nhiên có thể qua dữ liệu kiểm tra khi thi.

Nhưng với dữ liệu như sau thì sao? Hiển nhiên dữ liệu này vẫn phù hợp yêu cầu
đề bài:
\[
  (a+9999)^{9^{9^{9^{\cdot^{\cdot^{9}}}}}}.
\]
Không khó thấy rằng đa thức sau khi khai triển căn bản không thể lưu được,
và các hệ số phía sau cũng khó biểu diễn. Trước một đối tượng khổng lồ như
vậy, làm sao so sánh hai đa thức đều lớn đến thế có bằng nhau hay không? Rõ
ràng ngay cả tính toàn bộ biểu thức ra đã cực kỳ khó, chưa nói đến so sánh.

Thật ra cách giải rất đơn giản: thay một hoặc vài giá trị vào biến $a$, rồi
dựa vào phần dư của kết quả sau khi chia cho một số lớn để trực tiếp phán
định hai đa thức có bằng nhau hay không. Nếu không lấy dư, dữ liệu ``khổng
lồ'' ở trên vẫn không xử lý được. Việc thay một giá trị vào rồi tính kết quả
cũng có thể xem là áp dụng một hàm băm đặc biệt cho hai đa thức, sau đó so
sánh giá trị băm để so sánh biểu thức gốc. Dĩ nhiên, nếu chọn vài giá trị
phù hợp cho $a$, có thể chứng minh thuật toán này trở thành chắc chắn đúng.

\section{Ví dụ 6: Guess the Number}

\subsection{Mô tả bài toán}

Bài này được cải biên từ OIBH Reminiscence Programming Contest. Cho một số
$S$ có tối đa $500000$ chữ số. Hãy tìm $N$ sao cho $N^N=S$. Tuy nhiên, số
$S$ được cho có thể có một vài chữ số sai, nhưng không có chuyện thêm hoặc
bớt chữ số; khi đó cần in ra $-1$.

\subsection{Phân tích thuật toán}

Trước hết, với $S$ nhỏ có thể kiểm tra trực tiếp. Còn với $S>10^{10}$, ta có
thể xác định một giá trị $N$ khả dĩ, vì khi $N\ge 10$ thì
\[
  \frac{(N+1)^{N+1}}{N^N}>10,
\]
nên hai lũy thừa liên tiếp không thể có cùng số chữ số. Phần còn lại là kiểm
tra $N$ ấy có thỏa $N^N=S$ hay không.

Hiển nhiên tính logarit là không thể dựa vào, vì kiểu số thực không có độ
chính xác cao như vậy. Tương tự, tính số nguyên lớn cũng khó thực hiện; ngay
cả FFT cũng gặp thách thức về độ phức tạp thời gian. Liên hệ với các ví dụ
trước, ta có thể nghĩ đến việc tìm một hàm băm $f(x)$, rồi so sánh
$f(N^N)$ và $f(S)$ để xác định $N^N$ có bằng $S$ hay không. Đồng thời, hàm
$f(x)$ còn cần thỏa: nếu $x=ab$, ta có thể dễ dàng tính $f(x)$ từ $f(a)$ và
$f(b)$, nhờ đó so sánh trực tiếp mà không cần tính ra $N^N$.

Một hàm băm rất đơn giản nhưng hiệu quả là lấy dư:
\[
  f(x)=x\bmod p.
\]
Nếu $x=ab$ thì
\[
  f(x)=f(a)f(b)\bmod p.
\]

Với bài gốc, trong đó $S$ và $N^N$ nhiều nhất chỉ khác một chữ số, chỉ cần
$p$ không chia hết cho $10$ là có thể bảo đảm phán định chính xác.

Với phiên bản cải biên, nơi $S$ có thể sửa tùy ý nhiều chữ số, thì sao? Lấy
một vài giá trị $p$. Do giới hạn độ dài của $S$ chỉ là $500000$ chữ số, tích
các số nguyên tố nhỏ hơn $130000$ đã vượt quá phạm vi có thể chia hết bởi
$|S-N^N|$. Nói cách khác, nếu chọn ngẫu nhiên một $p$ trong phạm vi $500000$,
xác suất $p$ vừa khéo chia hết $|S-N^N|$ chỉ khoảng $1/4$. Nếu chọn nhiều
giá trị hơn, hoặc mở rộng phạm vi của $p$, ta gần như chắc chắn phán định
đúng.

\section*{Tổng kết}

Ngoài vai trò công cụ phụ trợ cho bảng băm, khi dùng độc lập, hàm băm có thể
phán định gần như chắc chắn hai dữ liệu có giống nhau hoặc tương đương hay
không. Bản chất của hàm băm là ``tóm lược'' thông tin có lượng dữ liệu lớn.
Vì vậy khi gặp bài toán cần thường xuyên so sánh hai dữ liệu có giống nhau
hay không, ta nên cân nhắc dùng hàm băm như một công cụ giải bài.

Đôi khi dùng hàm băm sẽ đánh đổi một phần tính đúng đắn, và thuật toán dùng
hàm băm nhìn qua cũng không đủ đẹp. Nhưng nếu dùng hàm băm linh hoạt, ta có
thể biến khó thành dễ, giải những bài vốn khó xử lý, và thường đạt hiệu quả
cao bằng lượng mã ngắn.

Dù có điều phải bỏ, cũng có điều nhận lại; chính vì có điều bỏ đi nên mới có
điều nhận được. Khi chọn thuật toán, ta luôn gặp kiểu mâu thuẫn như vậy: có
mất thì có được, có được thì cũng có mất. Thuật toán tốt nhất là thuật toán
phù hợp với đề bài và cũng phù hợp với thói quen tư duy của mình.

\section*{Tài liệu tham khảo}

\begin{itemize}
  \item Liu Rujia, Huang Liang, \textit{Nghệ thuật thuật toán và thi tin học},
  Nhà xuất bản Đại học Thanh Hoa.
  \item Chen Qifeng, \textit{NEW LCP}, bài tập đội tuyển tập huấn quốc gia
  năm 2006.
  \item \textit{Cấu trúc dữ liệu}, Nhà xuất bản Đại học Thanh Hoa.
\end{itemize}

\end{document}
